\documentclass[10pt,t,sans,mathsans,xcolor=dvipsnames]{beamer}

\usetheme{Antibes}
\usecolortheme{structure}
%\usepackage[utf8x]{inputenc}
%\usepackage{default}
%\usecolortheme{albatross}
%\usecolortheme{lily}
%\usecolortheme{sidebartab}
%\usecolortheme{crane}
%\usecolortheme{orchid}
%\usecolortheme{albatross}
%\usecolortheme{beetle}
%\usecolortheme{dove}
%\usecolortheme{fly}
%\usecolortheme{seagull}
%\usecolortheme{dolphin}
%\usecolortheme{rose}
\usepackage{graphics}
%\usepackage[pdftex]{graphicx}
%\usepackage{graphicx}
\usepackage{xcolor}
\usepackage{color}
%\usepackage{url}
\usepackage{hyperref}
%\usepackage[obeyspaces]{url}
\usepackage{amssymb,amsmath} % for the bold symbol command
\usepackage{booktabs} % toprule etc in tables
\usepackage{mathrsfs}
\usepackage{listings}
\lstset{ %
  backgroundcolor=\color{white},
  basicstyle=\footnotesize,
  breakatwhitespace=false,
  breaklines=true,
  captionpos=b,
  commentstyle=\color{green},
  escapeinside={\%*}{*)},
  extendedchars=true,
  frame=single,
  keywordstyle=\color{blue},
  language=bash,
  numbers=left,
  numbersep=5pt,
  numberstyle=\tiny\color{gray},
  rulecolor=\color{black},
  showspaces=false,
  showstringspaces=false,
  showtabs=false,
  stepnumber=2,
  stringstyle=\color{red},
  tabsize=2,
  title=\lstname,
  morekeywords={not,\},\{,preconditions,effects },
  deletekeywords={time}
}
%\usepackage{beamerthemebars}
%\usepackage{ragged2e}
%\usepackage{lipsum}

\setbeamercolor{structure}{fg=OliveGreen!50!black}
\renewcommand{\raggedright}{\leftskip=0pt \rightskip=0pt plus 0cm}
\logo{\includegraphics[scale=0.05]{../../../TWAssets/TW_Colour_Logos_trans_earth.png}}

\title{Introduction to version control with Git}
\titlegraphic{\includegraphics[scale=0.15]{../../../TWAssets/TW_Colour_Logos_trans_earth.png}}

\author{ Brain Leke Betetchouh \& 'Wole Solana }
\date{}
\begin{document}
\nocite*{}
\frame [c, plain]{\titlepage}
%-------------------------------Slide1----------------------------------------------------------------------------------------
\section{Introduction}
\begin{frame}
\frametitle{What is Git?}
\pause
\begin{itemize}[<+->]
\item Git is a \alert{version control} system (VCS) that allows you to keep track of changes made to a file, so you can recall specific versions of the file over time.
\item Other version control systems include centralized VCSs such as Subversion and distributed VCSs like Mercurial.
\end{itemize}
\end{frame}
%------------------------------Slide2-------------------------------------------------------------------------------------------
\section{The Basics}
\begin{frame}[fragile]
\frametitle{Setup}
\pause
\begin{itemize}[<+->]
\item Installing Git
\begin{lstlisting}
apt-get install git
\end{lstlisting}
\item Configuring your settings
\begin{lstlisting}
git config --global user.name "Your Name"
git config --global user.email your@emailaddress.com
\end{lstlisting}
\item Initializing a local repo
\begin{lstlisting}
git init
\end{lstlisting}
\end{itemize}
\end{frame}
%------------------------------Slide3-----------------------------------
\subsubsection*{Tracking Changes}
\begin{frame}[fragile, allowframebreaks]
\frametitle{Tracking changes}
\begin{itemize}
\item Checking the status of your repo
\begin{lstlisting}
git status
git diff
\end{lstlisting}
\item Undoing changes to files
\begin{itemize}
\item Single file:
\begin{lstlisting}
git checkout filename
\end{lstlisting}
\item Multiple files:
\begin{lstlisting}
git checkout .
\end{lstlisting}
\end{itemize}
\item[]
\item Staging and Unstaging files
\begin{itemize}
\item Single file:
\begin{lstlisting}
git add filename
\end{lstlisting}
\item Multiple files:
\begin{lstlisting}
git add .
\end{lstlisting}
\end{itemize}
\item[]
\item[]
\item[]
\item[]
\item Committing your files
\begin{itemize}
\item Single file:
\begin{lstlisting}
git commit filename
git commit -m "Commit message"
\end{lstlisting}
\item Multiple files:
\begin{lstlisting}
git commit .
\end{lstlisting}
\item There is a command to \textit{add} and \textit{commit} in one line. That is left as an exercise for the students.
\end{itemize}
\item Ignoring files
\begin{itemize}
\item Create a file `\texttt{.gitignore}` and add the filenames (or file patterns) you wish to ignore
\begin{lstlisting}
touch `.gitignore'
\end{lstlisting}
\end{itemize}
\item Amending a commit
\begin{lstlisting}
git commit --amend
\end{lstlisting}
\item Commit History
\begin{lstlisting}
git log
\end{lstlisting}
\end{itemize}
\end{frame}
%---------------------------------------------Slide4------------------------------------------------------------
\section{Working with remote repositories}
\begin{frame}[fragile]
\frametitle{Remote repositories}
\begin{itemize}[<+->]
\item Initializing from a remote repo
\begin{lstlisting}
git clone [remote_url]
\end{lstlisting}
\item Setting up a repo on GitHub, BitBucket, etc.
\item Adding a remote repo manually
\begin{lstlisting}
git remote add [remote_name] [remote_url]
\end{lstlisting}
\item Pulling from a remote
\begin{lstlisting}
git fetch [remote_name]
git pull [remote_name]
\end{lstlisting}
\item Pushing your work
\begin{lstlisting}
git push [remote_name]
\end{lstlisting}
\end{itemize}
\end{frame}
%---------------------------------------------Slide5----------------------------------------------------------------
\section{Just gitting started...}
\begin{frame}[fragile]%[allowframebreaks]
\frametitle{Just gitting started}
\pause
\begin{itemize}[<+->]
\item Branching
\item Merging and merge conflicts
\item Working with multiple remotes
\end{itemize}
\end{frame}
%---------------------------------------------Slide6------------------------------------------------------------------
\section{Questions and Exercises}
\begin{frame}[fragile]
\frametitle{Questions and Exercises}
\pause
\begin{itemize}[<+->]
\item Register yourself at \texttt{bitbucket.com} (if you don't have an account set up)
\item Email your user name to \texttt{wole@aims.ac.tz}
\item Using the commands introduced earlier, set up a git repo in a folder titled \texttt{TW\_AIMS\_Exercises}
\item Clone from this remote url:\\
\footnotesize{\texttt{https://bitbucket.org/wolesolana/aims-tz-student-exercises.git}}
\item Modify the file named \texttt{Participants.txt} by adding your name and email address in the format of the previous entries
\item Commit your changes with a good commit message
\item Push your changes
\end{itemize}
\end{frame}
%---------------------------------------------Slide7------------------------------------------------------------------
\section{Further Reading}
\begin{frame}
\frametitle{Further Reading}
\bibliographystyle{plain}
\bibliography{IntroToGit.bib}
\end{frame}
\end{document}
